% 10_call_to_action.tex — Call to Action
% ============================================================

\section{Call to Action}
\label{sec:action}

The infrastructure described in this paper is only as valuable as the
community that stress-tests, extends, and ultimately governs it.  We
outline concrete entry points for four audiences.

\paragraph{Professors.}
Co-steward the LIP-Arena evaluation tracks.  Validate your group's
methods on the four-scenario protocol (Ideal, Baseline, Oracle,
Full-Pipeline) and publish the full diagnostic vector---not just the
Scenario~I PSNR that flatters every solver equally.  Serve on the CISP
steward board and shape the metrics that will define the field for the
next decade.  The evaluation protocol is designed to be extensible: if
your modality is missing, proposing a new \texttt{OperatorGraph} template
earns co-authorship on the benchmark paper and a permanent seat at the
standards table.

\paragraph{PhD Students.}
Join the calibration sprints.  Pick one of the 64 registered modalities,
measure its mask-sensitivity spectrum on physical hardware, and
contribute a row to the correction table that currently spans 16
modalities.  Each new row is a quantified, reproducible contribution---a
first-author publication opportunity in a subfield that did not exist two
years ago.  The toolchain is open:

\begin{center}
\texttt{pip install pwm-eval \&\& pwm calibrate --modality <yours>}
\end{center}

\noindent
From installation to a publishable sensitivity curve takes one
afternoon; from curve to a peer-reviewed correction-factor paper takes
one semester.

\paragraph{Hobbyists.}
Participate in the weekly SolveEverything challenges.  Download a
measurement set, run your solver---classical, learned, or hybrid---and
submit a \texttt{RunBundle} for automated scoring against the blinded
ground truth.  The entire infrastructure is released under the MIT
license: zero institutional affiliation required, zero cost, zero
barrier to entry.  The leaderboard ranks results, not credentials.

\paragraph{Investors.}
The flywheel economics described in \cref{sec:flywheel} create three
revenue surfaces once the standard achieves critical mass:
\emph{calibration-as-a-service} (upload raw sensor data, receive a
validated operator), \emph{pay-per-reconstruction APIs} (submit a
measurement, receive a reconstruction with uncertainty bounds), and
\emph{imaging SLAs} (outcome-based contracts guaranteeing a minimum
PSNR/SSIM on the client's hardware).  The first evaluation standard
adopted by five or more laboratories locks in the market---not through
intellectual-property moats, but through the network effect of a shared
benchmark.  The competitive question is not \emph{whether} this standard
will emerge, but \emph{who} writes it.
